






























\subsection{特征值与特征向量}
接下来我们将着重讨论
\nearbylemma{lm:DiagIffBasisOfEigens} 的性质。

\begin{definition} \label{def:Eigen}
%<*df:Eigen>
变换 \( \map{t}{V}{V} \) 具有标量
\definend{特征值}\index{eigenvalue, eigenvector!of a transformation}%
\index{transformation!eigenvalue, eigenvector}
\( \lambda \)，是指存在非零的
\definend{特征向量} \( \vec{\zeta}\in V \)
使得
$
  t(\vec{\zeta})=\lambda\cdot\vec{\zeta}
$。
%</df:Eigen>
\end{definition}

\noindent ``Eigen'' 是德语词，意为``……所特有的''或``……独有的''。
一些作者将这些值与向量称为 \definend{特征}%
\index{characteristic!vectors, values} 值与特征向量。
没有作者称它们为``独有的''向量。

\begin{example}
投影映射
\begin{equation*}
  \colvec{x \\ y \\ z}
     \mapsunder{\pi}
  \colvec{x \\ y \\ 0}
   \qquad x,y,z\in\C
\end{equation*}
对任意形如
\begin{equation*}
   \colvec{x \\ y \\ 0}
\end{equation*}
（其中 \( x \) 与 \( y \) 是不全为零的标量）的特征向量，都有特征值 \( 1 \) 与之相伴。

相比之下，\( 2 \) 不是这个映射的特征值，
因为假设 $\pi$ 将一个向量加倍，
就会导出三个方程 $x=2x$、$y=2y$ 以及 $0=2z$，
于是没有非 $\zero$ 向量被加倍。
\end{example}

注意，
定义要求特征向量非 $\zero$。
一些作者允许 $\zero$ 作为 $\lambda$ 的特征向量，
只要同时还有与 $\lambda$ 相伴的
非 $\zero$ 向量即可。
关键在于
排除这样的平凡情形：$\lambda$ 满足
$t(\vec{v})=\lambda\vec{v}$ 仅对唯一的向量 $\vec{v}=\zero$ 成立。

另外，注意特征值 $\lambda$ 可以是~$0$。
问题在于 $\vec{\zeta}$ 是否等于 $\zero$。

\begin{example}  \label{ex:NoEigenOnTrivSp}
平凡空间 \( \set{\zero} \) 上仅有的变换是
%\begin{equation*}
$\zero\mapsto\zero$。
%\end{equation*}
这个映射没有特征值，因为不存在被映射到自身标量倍 $\lambda\cdot\vec{v}$ 的
非 \( \zero \) 向量 $\vec{v}$。
\end{example}

\begin{example} \label{ex:TransPolyOne}
考虑同态 \( \map{t}{\polyspace_1}{\polyspace_1} \)，
它由 \( c_0+c_1x\mapsto(c_0+c_1)+(c_0+c_1)x \) 给出。
虽然 $t$ 的陪域 $\polyspace_1$ 是二维的，但其
值域是一维的 $\rangespace{t}=\set{c+cx\suchthat c\in\C}$。
应用
\( t \) 作用到该值域中的向量只是对向量作伸缩：
\( c+cx\mapsto (2c)+(2c)x \)。
也就是说，\( t \) 有特征值 \( 2 \) 与形如 \( c+cx \)（其中 \( c\neq 0 \)）的
特征向量相伴。

这个映射也有特征值 \( 0 \)，与形如 \( c-cx \)（其中 \( c\neq 0 \)）的
特征向量相伴。
\end{example}

上面的定义是针对映射的。
我们也可以给出矩阵版本。

\begin{definition}  \label{df:EigenOfMatrix}
%<*df:EigenOfMatrix>
若方阵 \( T \) 满足 \( T\vec{\zeta}=\lambda\cdot\vec{\zeta} \)，
其中 \( \vec{\zeta} \) 是与之相伴的非零
\definend{特征向量}，则称 \( T \) 有标量
\definend{特征值}\index{eigenvalue, eigenvector!of a matrix}
\( \lambda \)。
%</df:EigenOfMatrix>
\end{definition}

把针对映射的定义扩展成针对矩阵的定义是很自然的，
但有一点必须小心。
映射的特征值也是表示该映射的矩阵的特征值，
所以相似的矩阵有相同的特征值。
然而，特征向量却可以
不同\Dash 相似的矩阵未必有
相同的特征向量。
下面的例子说明了这一点。

\begin{example}
这两个矩阵是相似的
\begin{equation*}
  T=
  \begin{mat}
    2  &0  \\
    0  &0
  \end{mat}
  \quad
  \hat{T}
  =
  \begin{mat}[r]
    4  &-2  \\
    4  &-2
  \end{mat}
\end{equation*}
因为对这个 $P$ 有 $\hat{T}=PTP^{-1}$。
\begin{equation*}
  P=
  \begin{mat}
    1  &1  \\
    1  &2
  \end{mat}
  \quad
  P^{-1}=
  \begin{mat}[r]
    2   &-1  \\
    -1  &1
  \end{mat}
\end{equation*}
矩阵 $T$
有两个特征值：$\lambda_1=2$ 和~$\lambda_2=0$。
第一个特征值与这个特征向量相伴。
\begin{equation*}
  T\vec{e}_1=
  \begin{mat}
    2  &0  \\
    0  &0
  \end{mat}
  \colvec{1 \\ 0}
  =\colvec{2 \\ 0}
  =2\vec{e}_1
\end{equation*}
假设 $T$ 是关于标准基表示变换 $\map{t}{\C^2}{\C^2}$
所得的矩阵。
那么这个变换 $t$ 的作用很简单。
\begin{equation*}
  \colvec{x \\ y}\mapsunder{t}\colvec{2x \\ 0}
\end{equation*}

当然，$\hat{T}$ 表示同一个变换，但
是关于不同的基~$B$。
我们可以求出这个基。
沿着箭头图从左下到左上，
\begin{equation*}
  \begin{CD}
    V_{\wrt{\stdbasis_2}}            @>t>T>        V_{\wrt{\stdbasis_2}}       \\
    @V{\scriptstyle\identity} VV              @V{\scriptstyle\identity} VV \\
    V_{\wrt{B}}                   @>t>\hat{T}>        V_{\wrt{B}}
  \end{CD}
\end{equation*}
可以看出 $P^{-1}=\rep{\identity}{B,\stdbasis_2}$。
根据映射的矩阵表示的定义，它的第一列是
$\rep{\identity(\vec{\beta}_1)}{\stdbasis_2}=\rep{\vec{\beta}_1}{\stdbasis_2}$。
关于标准基，任何向量都由其自身表示，
所以第一个基元素 $\vec{\beta}_1$ 就是 $P^{-1}$ 的第一列。
另一个基元素同理。
\begin{equation*}
  B=\sequence{\colvec[r]{2 \\ -1},
              \colvec[r]{-1 \\ 1}
              }
\end{equation*}
由于矩阵 $T$ 和~$\hat{T}$ 都表示变换~$t$，
它们都体现了作用 $t(\vec{e}_1)=2\vec{e}_1$。
\begin{align*}
  &\rep{t}{\stdbasis_2,\stdbasis_2}\cdot\rep{\vec{e}_1}{\stdbasis_2}
    =T\cdot\rep{\vec{e}_1}{\stdbasis_2}
    =2\cdot\rep{\vec{e}_1}{\stdbasis_2}                  \\
    &\rep{t}{B,B}\cdot\rep{\vec{e}_1}{B}
    =\hat{T}\cdot\rep{\vec{e}_1}{B}
    =2\cdot\rep{\vec{e}_1}{B}
\end{align*}
但是，虽然在这两个等式中特征值 $2$ 相同，
向量的表示却不同。
\begin{align*}
    T\cdot\rep{\vec{e}_1}{\stdbasis_2}
    =T\colvec{1 \\ 0}
    &=2\cdot\colvec{1 \\ 0}                \\
    \hat{T}\cdot\rep{\vec{e}_1}{B}
    =\hat{T}\cdot\colvec{1 \\ 1}
    &=2\cdot\colvec{1 \\ 1}
\end{align*}
也就是说，当表示变换的矩阵是
\( T=\rep{t}{\stdbasis_2,\stdbasis_2} \) 时，它``默认''
列向量是
关于 \( \stdbasis_2 \) 的表示。
而 \( \hat{T}=\rep{t}{B,B} \) 则``默认''列向量
是关于 \( B \) 的表示，
于是被加倍的那些列向量
是不同的。
\end{example}


% \begin{example}
% Consider again the transformation \( \map{t}{\polyspace_1}{\polyspace_1} \)
% % from \nearbyexample{ex:TransPolyOne}
% given by
% \( c_0+c_1x\mapsto (c_0+c_1)+(c_0+c_1)x \).
% One of its eigenvalues is \( 2 \), associated with the eigenvectors
% \( c+cx \) where \( c\neq 0 \).
% If we represent \( t \) with respect to \( B=\sequence{1+1x,1-1x} \)
% \begin{equation*}
%    T=\rep{t}{B,B}=
%    \begin{mat}[r]
%       2  &0  \\
%       0  &0
%    \end{mat}
% \end{equation*}
% then \( 2 \) is an eigenvalue of the matrix \( T \),
% associated with these eigenvectors.
% \begin{equation*}
%    \set{\colvec{c_0 \\ c_1}\suchthat \begin{mat}[r]
%                                          2  &0  \\
%                                          0  &0
%                                       \end{mat}\colvec{c_0 \\ c_1}
%                                       =\colvec{2c_0 \\ 2c_1}  }
%   =\set{\colvec{c_0 \\ 0}\suchthat c_0\in\C,\, c_0\neq 0 }
% \end{equation*}
% On the other hand, if we represent $t$ with respect to
% \( D=\sequence{2+1x,1+0x} \)
% \begin{equation*}
%    S=\rep{t}{D,D}=
%    \begin{mat}[r]
%       3  &1  \\
%      -3  &-1
%    \end{mat}
% \end{equation*}
% then the eigenvectors associated with the eigenvalue \( 2 \) are
% these.
% \begin{equation*}
%    \set{\colvec{c_0 \\ c_1}\suchthat \begin{mat}[r]
%                                          3  &1  \\
%                                         -3  &-1
%                                       \end{mat}\colvec{c_0 \\ c_1}
%                                       =\colvec{2c_0 \\ 2c_1}  }
%   =\set{\colvec{0 \\ c_1}\suchthat c_1\in\C,\, c_1\neq 0 }
% \end{equation*}
% \end{example}


我们接下来看到求
特征向量与特征值的基本工具。

\begin{example}  \label{ex:IntroCharEqn}
若
\begin{equation*}
  T=
  \begin{mat}[r]
     1    &2    &1    \\
     2    &0    &-2   \\
    -1    &2    &3
  \end{mat}
\end{equation*}
那么为求出使得
\( T\vec{\zeta}=x\vec{\zeta} \)（其中 \( \vec{\zeta} \) 为非零特征向量）的标量 \( x \)，
把所有项移到左边
\begin{equation*}
  \begin{mat}[r]
     1    &2    &1    \\
     2    &0    &-2   \\
    -1    &2    &3
  \end{mat}
  \colvec{z_1 \\ z_2 \\ z_3}
  -x\colvec{z_1 \\ z_2 \\ z_3}
  =\zero
\end{equation*}
然后因式分解，得到
\( (T-x I)\vec{\zeta}=\zero \)。
（注意是 $T-xI$。
式子 \( T-x \) 没有意义，
因为 \( T \) 是矩阵而 \( x \) 是标量。）
这个齐次线性方程组
\begin{equation*}
  \begin{mat}
   1-x           &2            &1            \\
     2           &0-x          &-2           \\
    -1           &2            &3-x
  \end{mat}
  \colvec{z_1 \\ z_2 \\ z_3}
  =
  \colvec[r]{0 \\ 0 \\ 0}
\end{equation*}
有非零解 $\vec{z}$ 当且仅当该
矩阵是奇异的。
我们可以确定何时会发生这种情况。
\begin{align*}
  0
  &=\deter{T-x I}                                               \\
  &=\begin{vmatrix}
     1-x          &2            &1            \\
     2           &0-x          &-2           \\
    -1           &2            &3-x
   \end{vmatrix}                                       \\
  &=x^3-4x^2+4x  \\
  &=x(x-2)^2
\end{align*}
特征值是 \( \lambda_1=0 \) 和 \( \lambda_2=2 \)。
为求相伴的特征向量，将每个特征值代入。
代入 $\lambda_1=0$ 给出
\begin{equation*}
  \begin{mat}
     1-0         &2            &1            \\
     2           &0-0          &-2           \\
    -1           &2            &3-0
  \end{mat}
  \colvec{z_1 \\ z_2 \\ z_3}
  =
  \colvec[r]{0 \\ 0 \\ 0}
  \quad\Longrightarrow\quad
  \colvec{z_1 \\ z_2 \\ z_3}
  =
  \colvec[r]{a \\ -a \\ a}
\end{equation*}
其中 \( a\neq 0 \)
（\( a \) 必须非 $0$，因为按定义特征向量
是非 \( \zero \) 的）。
代入 $\lambda_2=2$ 给出
\begin{equation*}
  \begin{mat}
     1-2         &2            &1            \\
     2           &0-2          &-2           \\
    -1           &2            &3-2
  \end{mat}
  \colvec{z_1 \\ z_2 \\ z_3}
  =
  \colvec[r]{0 \\ 0 \\ 0}
  \quad\Longrightarrow\quad
  \colvec{z_1 \\ z_2 \\ z_3}
  =
  \colvec{b \\ 0 \\ b}
\end{equation*}
其中 \( b\neq 0 \)。
\end{example}

\begin{example} \label{ex:AnotherCharPoly}
若
\begin{equation*}
  S=
  \begin{mat}[r]
    \pi      &1      \\
    0        &3
  \end{mat}
\end{equation*}
（这里 \( \pi \) 不是投影映射，而是数
\( 3.14\ldots \)），那么
\begin{equation*}
    \begin{vmat}
      \pi-x &1         \\
      0     &3-x
    \end{vmat}
  =
  (x-\pi)(x-3)
\end{equation*}
所以 \( S \) 有特征值 \( \lambda_1=\pi \) 和 \( \lambda_2=3 \)。
为求相伴的特征向量，先将 $\lambda_1$ 代入 $x$
\begin{equation*}
  \begin{mat}
    \pi-\pi     &1         \\
    0           &3-\pi
  \end{mat}
  \colvec{z_1 \\ z_2}
  =
  \colvec[r]{0 \\ 0}
  \qquad\Longrightarrow\qquad
  \colvec{z_1 \\ z_2}
  =
  \colvec{a \\ 0}
\end{equation*}
其中标量 \( a\neq 0 \)。
再将 $\lambda_2$ 代入
\begin{equation*}
  \begin{mat}
    \pi-3       &1         \\
    0           &3-3
  \end{mat}
  \colvec{z_1 \\ z_2}
  =
  \colvec[r]{0 \\ 0}
  \qquad\Longrightarrow\qquad
  \colvec{z_1 \\ z_2}
  =
  \colvec{-b/(\pi-3) \\ b}
\end{equation*}
其中 \( b\neq 0 \)。
\end{example}

\begin{definition} \label{df:CharacteristicPoly}
%<*df:CharacteristicPoly>
\( \nbyn{n} \) 方阵
\definend{特征多项式}\index{characteristic!polynomial}%
\index{matrix!characteristic polynomial}
\( T \) 是行列式 \( \deter{T-x I} \)，其中 \( x \) 是变量。
\definend{特征方程}\index{characteristic!equation}%
\index{matrix!characteristic polynomial}
是 $\deter{T-xI}=0$。
变换的
\definend{特征多项式}
\( t \) 是它的任意矩阵表示 \( \rep{t}{B,B} \)
的特征多项式。\index{transformation!characteristic polynomial}
%</df:CharacteristicPoly>
\end{definition}

\noindent
%<*df:CharacteristicPolyExer>
\( \nbyn{n} \) 矩阵
或变换 \( \map{t}{\C^n}{\C^n} \) 的特征多项式的次数为~\( n \)。
\nearbyexercise{exer:CharPolyTransWellDefed} 验证了变换的
特征多项式是
良定义的，也就是说，无论用哪个基作表示，特征多项式都是
相同的。
%</df:CharacteristicPolyExer>

\begin{lemma} \label{le:MapNonTrivSpHasEigen}
%<*lm:MapNonTrivSpHasEigen>
非平凡向量空间上的线性变换至少有
一个特征值。
%</lm:MapNonTrivSpHasEigen>
\end{lemma}

\begin{proof}
%<*pf:MapNonTrivSpHasEigen>
特征多项式的任何根都是特征值。
在复数域上，任何次数为一或更高的多项式
都有根。
%</pf:MapNonTrivSpHasEigen>
\end{proof}

\begin{remark}
这个结果正是本章
使用复数标量的原因。
如果坚持使用实数标量，那么就会有些特征多项式，
比如 $x^2+1$，不能分解。
\end{remark}

% Notice the familiar form of the sets of eigenvectors in the above examples.

\begin{definition} \label{df:Eigenspace}
%<*df:Eigenspace>
变换~$t$ 关于
特征值~$\lambda$
的 \definend{特征子空间}\index{eigenspace}\index{transformation!eigenspace}
是
$
  V_{\lambda}=\set{\vec{\zeta}\suchthat t(\vec{\zeta}\,)=\lambda\vec{\zeta}\,}
              % \union\set{\zero\,}
$。
矩阵的特征子空间与之类似。
%</df:Eigenspace>
\end{definition}

% \begin{remark}
% In \nearbydefinition{def:Eigen} we specified that eigenvectors must be
% non-$\zero$.
% So, strictly speaking, the eigenspace for $\lambda$ contains the set of
% associated eigenvectors along with $\zero$.
% We need the zero vector for the next result.
% \end{remark}

\begin{lemma}  \label{le:EigSpaceIsSubSp}
%<*lm:EigSpaceIsSubSp>
特征子空间是一个子空间。
它是非平凡的子空间。
%</lm:EigSpaceIsSubSp>
\end{lemma}

\begin{proof}
%<*pf:EigSpaceIsSubSp>
首先注意，
$V_{\lambda}$ 非空；它包含零
向量，因为 $t(\zero)=\zero$，而后者等于
$\lambda\cdot \zero$。
为证明特征子空间是子空间，
剩下只需检查这个集合对线性组合的封闭性。
取 \( \vec{\zeta}_1,\ldots,\vec{\zeta}_n\in V_{\lambda} \)，那么
\begin{align*}
  t(\lincombo{c}{\vec{\zeta}})
  &=c_1t(\vec{\zeta}_1)+\dots+c_nt(\vec{\zeta}_n)               \\
  &=c_1\lambda\vec{\zeta}_1+\dots+c_n\lambda\vec{\zeta}_n          \\
  &=\lambda(c_1\vec{\zeta}_1+\dots+c_n\vec{\zeta}_n)
\end{align*}
即该组合也是 $V_{\lambda}$ 的元素。
% (despite that the zero vector isn't an eigenvector,
% the second equality holds even if some \( \vec{\zeta}_i \) is \( \zero \) since
% \( t(\zero)=\lambda\cdot\zero=\zero \)).

空间~$V_\lambda$ 包含的不只是零向量，
因为按定义，只有当
$t(\vec{\zeta}\,)=\lambda\vec{\zeta}$ 有 $\zero$ 之外的
$\vec{\zeta}$ 解时，$\lambda$ 才是特征值。
%</pf:EigSpaceIsSubSp>
\end{proof}

\begin{example}   \label{ex:AlgMultDifferentThanGeoMult}
这些是与 \nearbyexample{ex:IntroCharEqn} 的特征值 \( 0 \)
和 \( 2 \) 相伴的特征子空间。
\begin{equation*}
  V_0=\set{\colvec[r]{a \\ -a \\ a}\suchthat a\in\C},
  \qquad
  V_2=\set{\colvec{b \\ 0 \\ b}\suchthat b\in\C}.
\end{equation*}
\end{example}

\begin{example} \label{ex:AnotherCharPolyReprise}
这些是 \nearbyexample{ex:AnotherCharPoly} 的特征值 \( \pi \)
和~\( 3 \)
对应的特征子空间。
\begin{equation*}
  V_{\pi}=\set{\colvec{a \\ 0}\suchthat a\in\C}
  \qquad
  V_3=\set{\colvec{-b/(\pi-3) \\ b}\suchthat b\in\C}
\end{equation*}
\end{example}

\nearbyexample{ex:IntroCharEqn} 中的
特征方程
是 \( 0=x(x-2)^2 \)，所以在某种意义上
\( 2 \) 两次成为特征值。
但特征向量的数目并没有翻倍，因为相伴特征子空间 $V_2$
的维数是一，不是二。
下一个例子则是一个数成为特征方程的二重根、
而相伴特征子空间
的维数是二的情形。

\begin{example}  \label{ex:EigenvaluesProjection}
关于标准基，这个矩阵
\begin{equation*}
  \begin{mat}[r]
     1  &0  &0  \\
     0  &1  &0  \\
     0  &0  &0
  \end{mat}
\end{equation*}
表示投影。
\begin{equation*}
  \colvec{x \\ y \\ z}
     \mapsunder{\pi}
  \colvec{x \\ y \\ 0}
   \qquad x,y,z\in\C
\end{equation*}
它的特征方程
\begin{align*}
  0
  &=\deter{T-x I}                                     \\
  &=\begin{vmatrix}
     1-x          &0            &0            \\
     0           &1-x          &0           \\
     0           &0            &0-x
   \end{vmatrix}                                       \\
  &=(1-x)^2(0-x)
\end{align*}
有二重根~$x=1$ 以及单根~$x=0$。
它关于特征值 \( 1 \) 的特征子空间
与关于特征值 \( 0 \) 的特征子空间
都容易求出。
\begin{equation*}
   V_1=\set{\colvec{c_1 \\ c_2 \\ 0}\suchthat c_1,c_2\in\C}
   \qquad
   V_0=\set{\colvec{0 \\ 0 \\ c_3}\suchthat c_3\in\C}
\end{equation*}
注意 $V_1$ 的维数是二。
\end{example}

\begin{definition}
当特征多项式分解为
$(x-\lambda_1)^{m_1}\cdots (x-\lambda_k)^{m_k}$ 时，
就称特征值 $\lambda_i$ 有
\definend{代数重数}~$m_i$\index{multiplicity!algebraic}%
\index{algebraic multiplicity}。
它的
\definend{几何重数}\index{multiplicity!geometric}%
\index{geometric multiplicity}
是相伴特征子空间~$V_{\lambda_i}$ 的维数。
\end{definition}

在 \nearbyexample{ex:EigenvaluesProjection} 中，有两个特征值。
对于 $\lambda_1=1$，
代数重数与几何重数都是~$2$。
对于 $\lambda_2=0$，
代数重数与几何重数都是~$1$。
相比之下，\nearbyexample{ex:AlgMultDifferentThanGeoMult}
表明特征值 $\lambda=2$ 的代数重数是~$2$，
但几何重数是~$1$。
对每个变换，由 \nearbylemma{le:EigSpaceIsSubSp} 知每个特征值的几何
重数都大于或等于~$1$。
（而且，特征值的几何重数必定
小于或等于
其代数重数，不过这个证明超出了本书的范围。）

由 \nearbylemma{le:EigSpaceIsSubSp}
知，若两个特征向量 $\vec{v}_1$ 与 $\vec{v}_2$
与同一个特征值相伴，那么这两个向量的线性组合也是特征向量，
且与同一个特征值相伴。
作为说明，参考前面的例题，
$V_1$ 中两个成员之和
\begin{equation*}
  \colvec[r]{1 \\ 0 \\ 0}+\colvec[r]{0 \\ 1 \\ 0}
\end{equation*}
给出 $V_1$ 的另一个成员。

下一个结果处理向量来自
不同特征子空间的情形。

\begin{theorem}  \label{th:DistEValueGivesLIEvecs}
%<*th:DistEValueGivesLIEvecs>
对于映射或矩阵的任意一组互不相同的特征值，一组与之相伴的
特征向量——每个特征值取一个——是线性无关的。
%</th:DistEValueGivesLIEvecs>
\end{theorem}

\begin{proof}
%<*pf:DistEValueGivesLIEvecs0>
我们将对特征值的个数使用归纳法。
基础步骤是特征值个数为零的情形。
此时相伴向量的集合是空集，
从而线性无关。
% If there is
% only one eigenvalue then the set of associated eigenvectors is
% a singleton set with a non-$\zero$ member,
% and so is linearly independent.
%</pf:DistEValueGivesLIEvecs0>

%<*pf:DistEValueGivesLIEvecs1>
对于归纳步骤，假设结论对任意由
\( k\geq 0 \) 个互不相同特征值组成的集合成立。
考虑互不相同的特征值
\( \lambda_1,\dots,\lambda_{k+1} \)，
并设 \( \vec{v}_1,\dots,\vec{v}_{k+1} \)
是与之相伴的特征向量。
假设
$\zero=c_1\vec{v}_1+\dots+c_k\vec{v}_k+c_{k+1}\vec{v}_{k+1}$。
由它导出两个方程：第一个是在两边乘以 \( \lambda_{k+1} \)
得到的
$\zero=c_1\lambda_{k+1}\vec{v}_1+\dots+c_{k+1}\lambda_{k+1}\vec{v}_{k+1}$，
第二个是把映射作用到两边得到的
$\zero=c_1t(\vec{v}_1)+\dots+c_{k+1}t(\vec{v}_{k+1})
  =c_1\lambda_1\vec{v}_1+\dots+c_{k+1}\lambda_{k+1}\vec{v}_{k+1}$
（用矩阵来做结果相同）。
用第一个减去第二个。
\begin{equation*}
  \zero=
  c_1(\lambda_{k+1}-\lambda_1)\vec{v}_1+\dots
  +c_k(\lambda_{k+1}-\lambda_k)\vec{v}_k
      +c_{k+1}(\lambda_{k+1}-\lambda_{k+1})\vec{v}_{k+1}
\end{equation*}
$\vec{v}_{k+1}$ 项消失。
于是归纳假设给出
\( c_1(\lambda_{k+1}-\lambda_1)=0 \), \ldots, \( c_k(\lambda_{k+1}-\lambda_k)=0 \)。
特征值互不相同，所以
系数 \( c_1,\,\dots,\,c_k \) 全都是~\( 0 \)。
这样一来
只剩方程 \( \zero=c_{k+1}\vec{v}_{k+1} \)，
所以 \( c_{k+1} \) 也是~\( 0 \)。
%</pf:DistEValueGivesLIEvecs1>
\end{proof}

\begin{example}
这个矩阵
\begin{equation*}
     \begin{mat}[r]
        2   &-2   &2   \\
        0   &1    &1   \\
       -4   &8    &3
     \end{mat}
\end{equation*}
的特征值是互不相同的：\( \lambda_1=1 \)，\( \lambda_2=2 \)，以及~\( \lambda_3=3 \)。
一组相伴的特征向量
\begin{equation*}
  \set{
       \colvec[r]{2 \\ 1 \\ 0},
       \colvec[r]{9 \\ 4 \\ 4},
       \colvec[r]{2 \\ 1 \\ 2}  }
\end{equation*}
是线性无关的。
\end{example}

\begin{corollary} \label{co:DistinctEivenvaluesImpliesDiagonal}
%<*co:DistinctEivenvaluesImpliesDiagonal>
有 \( n \) 个互不相同特征值的 \( \nbyn{n} \) 矩阵是可对角化的。
%</co:DistinctEivenvaluesImpliesDiagonal>
\end{corollary}

\begin{proof}
%<*pf:DistinctEivenvaluesImpliesDiagonal>
构造一个由特征向量组成的基。
应用 \nearbylemma{lm:DiagIffBasisOfEigens}。
%</pf:DistinctEivenvaluesImpliesDiagonal>
\end{proof}

\begin{example} \label{ex:SummaryOfDiagonalizable}
在前面的例题中我们证明了矩阵
\begin{equation*}
     T=\begin{mat}[r]
        2   &-2   &2   \\
        0   &1    &1   \\
       -4   &8    &3
     \end{mat}
\end{equation*}
是可对角化的，其特征值为
$\lambda_1=1$，$\lambda_2=2$，以及~$\lambda_3=3$。
这些是相伴的特征向量，它们组成一个基~$B$。
\begin{equation*}
       \vec{\beta}_1=\colvec[r]{2 \\ 1 \\ 0}\quad
       \vec{\beta}_2=\colvec[r]{9 \\ 4 \\ 4}\quad
       \vec{\beta}_3=\colvec[r]{2 \\ 1 \\ 2}
\end{equation*}
箭头图
\begin{equation*}
  \begin{CD}
    V_{\wrt{\stdbasis_3}}            @>t>T>        V_{\wrt{\stdbasis_3}}       \\
    @V{\scriptstyle\identity} VV              @V{\scriptstyle\identity} VV \\
    V_{\wrt{B}}                   @>t>D>        V_{\wrt{B}}
  \end{CD}
\end{equation*}
给出如下结果，其中 $P=\rep{\identity}{\stdbasis_3,B}$。
\begin{align*}
  D &= PTP^{-1}  \\
  \begin{mat}
    1  &0  &0 \\
    0  &2  &0 \\
    0  &0  &3
  \end{mat}
  &=
    \begin{mat}
      -2  &5  &-1/2  \\
       1  &-2  &0    \\
      -2  &4  &1/2
    \end{mat}
     \begin{mat}[r]
        2   &-2   &2   \\
        0   &1    &1   \\
       -4   &8    &3
     \end{mat}
    \begin{mat}
      2  &9  &2  \\
      1  &4  &1  \\
      0  &4  &2
    \end{mat}
\end{align*}
% sage: T = matrix(QQ, [[2,-2,2], [0,1,1], [-4,8,3]])
% sage: T
% [ 2 -2  2]
% [ 0  1  1]
% [-4  8  3]
% sage: T.eigenvectors_right()
% [(3, [
%   (1, 1/2, 1)
%   ], 1), (2, [
%   (1, 4/9, 4/9)
%   ], 1), (1, [
%   (1, 1/2, 0)
%   ], 1)]
% sage: P_inv = matrix(QQ, [[2,9,2], [1,4,1], [0,4,2]])
% sage: P_inv.inverse()
% [  -2    5 -1/2]
% [   1   -2    0]
% [  -2    4  1/2]
图的下半部分有
$t(\vec{\beta}_1)=1\cdot\vec{\beta}_1$，
还有 $t(\vec{\beta}_2)=2\cdot\vec{\beta}_2$，
以及 $t(\vec{\beta}_3)=3\cdot\vec{\beta}_3$。
改写即得：映射 $t-1$ 把 $\vec{\beta}_1$ 送到~$\zero$，
$t-2$ 把 $\vec{\beta}_2$ 送到~$\zero$，
而 $(t-3)(\vec{\beta}_3)=\zero$。
也就是说，在 $B$ 上的作用如下。
\begin{equation*}
  \vec{\beta}_1\mapsunder{t-1}\zero  \qquad
  \vec{\beta}_2\mapsunder{t-2}\zero  \qquad
  \vec{\beta}_3\mapsunder{t-3}\zero
\end{equation*}
转来看那个箭头图中的表示，最上面一行说的是：
表示 $t-1$ 关于
$\stdbasis_3,\stdbasis_3$ 的矩阵 $T-(1\cdot I)$，乘以
$\vec{\beta}_1$ 关于~$\stdbasis_3$ 的表示，将得到零。
\begin{equation*}
  \begin{mat}
    1  &-2  &2  \\
    0  &0   &1  \\
   -4  &8   &2  \\
  \end{mat}
  \colvec{2 \\ 1 \\ 0}
  =
  \colvec{0  \\ 0 \\ 0}
\end{equation*}
类似地，表示 $t-2$ 关于
$\stdbasis_3,\stdbasis_3$ 的矩阵 $T-(2\cdot I)$，乘以
$\vec{\beta}_2$ 关于~$\stdbasis_3$ 的表示，得到零。
\begin{equation*}
  \begin{mat}
    0  &-2  &2  \\
    0  &01   &1  \\
   -4  &8   &1  \\
  \end{mat}
  \colvec{9 \\ 4 \\ 4}
  =
  \colvec{0  \\ 0 \\ 0}
\end{equation*}
当然，矩阵 $T-(3\cdot I)$ 乘以
$\vec{\beta}_3$ 关于~$\stdbasis_3$ 的
表示也得到零。
\end{example}

总而言之，本节观察到一些
矩阵相似于对角矩阵。
特征值的想法正是作为那个对角矩阵的对角元而产生的，
尽管定义的适用范围比仅限于
可对角化矩阵更广。
为求特征值，我们定义了特征方程，
这引出了最后一个结果：一个
可对角化的判别法。
（虽然它对理论很有用，但请注意，
在实际应用中以这种方式求特征值
通常并不现实；一方面矩阵可能很大，
而求高次多项式的根是困难的。）

下一节我们将研究不能对角化的矩阵。


